\chapter{Benchmarks and microbenchmarks}
\label{ch:bench}

\begin{aside}
``It's fast'' is a feeling. Numbers are how we keep that
feeling honest. \maya{} ships microbenchmarks for the hot paths
so a refactor that quietly halves throughput shows up
immediately.
\end{aside}

\section{What to benchmark}

The hot paths:

\begin{itemize}[leftmargin=*]
  \item Canvas diff (SIMD).
  \item ANSI emit.
  \item Layout for a tree of $N$ elements.
  \item Signal-graph propagation.
  \item Grapheme cluster iteration.
  \item Width lookup.
\end{itemize}

The not-so-hot paths:

\begin{itemize}[leftmargin=*]
  \item App startup time.
  \item First-frame latency.
\end{itemize}

\section{Microbench framework}

\maya{} uses \code{nanobench} (header-only, single file):

\begin{lstlisting}
ankerl::nanobench::Bench()
    .run("diff 80x24 unchanged", [&] {
        diff_no_op(front, back);
    });
\end{lstlisting}

The framework auto-tunes iteration counts, reports
mean/median/stddev, and outputs in a parseable format.

\section{Comparing variants}

Run two variants and compare:

\begin{lstlisting}
Bench b;
b.run("scalar diff", [&] { scalar_diff(...); });
b.run("simd diff",   [&] { simd_diff(...); });
\end{lstlisting}

Output shows nanoseconds per op and ops per second.

\section{Macrobench}

A full app's frame loop is the macrobench. Drive a simulated
input stream and measure frames per second:

\begin{lstlisting}
TEST_CASE("60Hz typing") {
    auto app = make_app();
    for (int i = 0; i < 600; ++i) {
        app.send_key('a');
        app.render_frame();
    }
    // assert no frame > 16ms
}
\end{lstlisting}

This catches whole-pipeline regressions.

\section{Calibration}

Microbenchmarks need a baseline. \maya{} measures the
overhead of an empty loop and subtracts. The framework does
this automatically.

\section{Reference machine}

\maya{}'s CI runs benchmarks on a fixed machine. Numbers
are stored; a 10\% regression fails the run. The numbers do
not need to be portable --- relative regression is what
matters.

\section{Profiling vs benchmarking}

A benchmark measures throughput. A profile shows where time
goes. Both are needed:

\begin{itemize}[leftmargin=*]
  \item Bench first: ``the diff is 30\% slower''.
  \item Profile next: ``the wide-cell branch is mispredicting''.
  \item Fix and re-bench.
\end{itemize}

\section{Memory benchmarks}

Throughput isn't the only metric. Working-set size matters
for cache behaviour:

\begin{itemize}[leftmargin=*]
  \item Canvas: $w \times h \times 8$ bytes.
  \item Element tree: typically under 100 KB.
  \item Signal storage: under 10 KB for small apps.
\end{itemize}

\maya{} prints these on \code{about()}.

\section{Pitfalls}

\begin{itemize}[leftmargin=*]
  \item \textbf{Dead-code elimination.} The compiler may
    elide a benchmark whose result is unused. Mark with
    \code{doNotOptimizeAway}.
  \item \textbf{Cold cache.} First iteration is slow. Use
    median or warmup.
  \item \textbf{Frequency scaling.} CPU may downclock on
    battery. Pin the test machine.
  \item \textbf{Background noise.} Measure on idle systems
    or take the minimum.
\end{itemize}

\section{Reporting}

\maya{}'s CI posts benchmark deltas as PR comments. The
table shows ``before / after / delta'' for each named
benchmark. Reviewer sees the impact at a glance.

\section{Recap}

\begin{recap}
\begin{itemize}[leftmargin=*]
  \item Microbench the hot paths.
  \item Macrobench the full pipeline.
  \item Compare variants, not absolutes.
  \item Profile after benchmarking.
  \item Mind dead-code elimination and frequency scaling.
\end{itemize}
\end{recap}

\section*{Workshop}

\begin{enumerate}[leftmargin=*]
  \item Write a microbench for your most expensive widget.
    Bench it before and after a simple optimisation.
  \item Add a 60-Hz typing macrobench to your test suite.
  \item Compare scalar and SIMD diff on your CPU. Record the
    ratio.
\end{enumerate}
